\section{Edit-history replay}\label{sec:replay}

\subsection{Protocol and collection}
\code{experiments/revision_bundles/protocol.json} fixed the sampling frames,
the 2024 window, the capture schedules, the methods and the ground-truth
definitions, and was committed to the repository before any revision metadata
was requested; the collection manifest records its SHA-256
(\code{79c8be55...}), which \code{derive.py} checks. \code{collect.py} then queried the
English Wikipedia Action API serially, at least 2\,s apart, with
\code{maxlag=5}, retries after HTTP 429 and a descriptive User-Agent. It was
interrupted three times when its worker restarted (after 287, 318 and 359
saved responses); each time \code{--resume} replayed the saved responses
through the same request code to rebuild URLs and continuation tokens, failed
on any divergence, and sent only missing requests. The collection holds
\nRevisionRequests{} responses, retrieved between 05:55 and 10:48 UTC on
2026-11-05. The vital frame is the link table of \emph{Wikipedia:Vital
articles/Level~3} at revision 1382205991 (999 titles); the random frame is
MediaWiki's \code{page_random}. Only revision metadata (IDs, parents,
timestamps, SHA-1 values, sizes) was requested.

Of 300 sampled pages, 16 are excluded as the protocol requires: \nExcludedHidden{}
vital pages contain revisions whose SHA-1 is hidden, and \nExcludedAbsent{}
random pages have no visible revision by the end of the window. The
\nPages{} remaining pages (\nVitalPages{} vital, \nRandomPages{} random) give
\nStates{} states of positive duration (\sStatesVital{} vital,
\sStatesRandom{} random). \sPagesReturnVital{} vital and \sPagesReturnRandom{}
random pages return to an earlier SHA-1 at least once, in \sReturnEvents{}
return events, with median departures of \sMedianVital{}\,s (vital) and
\sMedianRandom{}\,s (random). The direction of selection bias from the
exclusions is unknown; the shares and page-bootstrap intervals below describe
the retained histories and do not correct for it.

The replay reconstructs an edit-only timeline from the visible revision
metadata, completely paginated, with revisions that share a timestamp
collapsed by the preregistered one-second convention. Administrative deletion,
undeletion and other availability events were not collected, so the replay
evaluates certification against the reconstructed edit history, not against
what an unqualified current endpoint served over time. Captures are replayed
against the frozen histories; no request was sent to Wikipedia at the
simulated instants.

\subsection{Equal hashes that hide a change}
Capture instants are hourly over the window; a pair counts when the page
exists at both instants. A pair \emph{hides a change} when its endpoint hashes
are equal but a state with another SHA-1 was current for positive time in
between. Table~\ref{tab:hidden} gives the preregistered measurement behind
Figure~2a. Equal revision IDs hid no change at any gap.

\begin{table*}[t]
\caption{Equal-hash revisit pairs that hid a change, with 95\% bootstrap intervals over pages (2{,}000 resamples, seed 7).}
\label{tab:hidden}
\small
\begin{tabular}{@{}llrrrlr@{}}
\toprule
Stratum & Gap & Pairs & Equal hash & Hidden & Share, \% (95\% CI) & Pages \\
\midrule
vital & 1\,min & 1,212,192 & 1,211,925 & 4 & ${<}0.01$ (${<}0.01$--${<}0.01$) & 4 of 138 \\
vital & 1\,h & 1,212,054 & 1,202,086 & 999 & 0.08 (0.07--0.10) & 118 of 138 \\
vital & 1\,d & 1,208,880 & 1,060,423 & 27,445 & 2.59 (2.14--3.09) & 127 of 138 \\
vital & 7\,d & 1,189,008 & 611,897 & 81,712 & 13.35 (11.04--15.80) & 129 of 138 \\
random & 1\,min & 1,263,461 & 1,263,409 & 1 & ${<}0.01$ (${<}0.01$--${<}0.01$) & 1 of 146 \\
random & 1\,h & 1,263,315 & 1,262,110 & 41 & ${<}0.01$ (${<}0.01$--${<}0.01$) & 12 of 146 \\
random & 1\,d & 1,259,957 & 1,246,560 & 859 & 0.07 (0.02--0.15) & 12 of 146 \\
random & 7\,d & 1,238,933 & 1,172,040 & 3,298 & 0.28 (0.10--0.52) & 12 of 146 \\
\bottomrule
\end{tabular}

\end{table*}

\subsection{Bundle certification}
The request gap is the time between consecutive request starts; every request
window lasts 1\,s, so at a 1\,s gap consecutive windows touch, and the minute
and hour gaps separate them further. Two pages is the smallest bundle, and four
and eight double it. The grid varies bundle size and timing under control; it
does not model the size or pacing of real crawls. Each stratum, bundle size ($n=2,4,8$) and gap uses 400 bundles drawn from the
pages that exist throughout the window (\sPoolVital{} vital and \sPoolRandom{}
random; the other \sPoolRandomLate{} random pages first appear during 2024) and
10 seeded start times: 4{,}000 runs per method and
condition, 12{,}000 per stratum and gap. The preregistered checks are G1, that
every page had its observed content at the claimed common time (for the exact
test, at the single read's instant), and G2, that no twice-read page changed
between its reads. Methods marked $^{\ast}$ are preregistered.

\begin{table*}[t]
\caption{Preregistered methods: certified runs and runs violating G1 and G2. Rows ``all $n$'' total 12{,}000 runs over $n=2,4,8$; the last three rows split vital pages at 1\,h by bundle size (4{,}000 runs each).}
\label{tab:prereg}
\small
\begin{tabular}{@{}lrrrrrrrrr@{}}
\toprule
Condition & \multicolumn{3}{c}{Content hash, $2n$, half-open} & \multicolumn{3}{c}{Revision ID, $2n$, half-open} & \multicolumn{3}{c}{Revision ID, exact, $2n{-}1$} \\ \cmidrule(lr){2-4}\cmidrule(lr){5-7}\cmidrule(l){8-10}
 & Certified & G1 & G2 & Certified & G1 & G2 & Certified & G1 & G2 \\
\midrule
vital, 1\,s, all $n$ & 0 & 0 & 0 & 0 & 0 & 0 & 11,997 & 0 & 0 \\
vital, 1\,min, all $n$ & 11,929 & 0 & 5 & 11,924 & 0 & 0 & 11,941 & 0 & 0 \\
vital, 1\,h, all $n$ & 10,195 & 8 & 241 & 9,954 & 0 & 0 & 10,279 & 0 & 0 \\
random, 1\,s, all $n$ & 0 & 0 & 0 & 0 & 0 & 0 & 12,000 & 0 & 0 \\
random, 1\,min, all $n$ & 11,997 & 0 & 0 & 11,997 & 0 & 0 & 11,997 & 0 & 0 \\
random, 1\,h, all $n$ & 11,852 & 0 & 4 & 11,848 & 0 & 0 & 11,869 & 0 & 0 \\
\midrule
vital, 1\,h, $n=2$ & 3,882 & 1 & 14 & 3,868 & 0 & 0 & 3,930 & 0 & 0 \\
vital, 1\,h, $n=4$ & 3,606 & 0 & 59 & 3,547 & 0 & 0 & 3,659 & 0 & 0 \\
vital, 1\,h, $n=8$ & 2,707 & 7 & 168 & 2,539 & 0 & 0 & 2,690 & 0 & 0 \\
\bottomrule
\end{tabular}

\end{table*}

Two observations follow from the preregistered counts
(Table~\ref{tab:prereg}). In every condition the surplus of the hash double collect over the
revision-ID double collect, in certified runs, equals its number of G2
violations; it is positive in \sPreregMore{} of the \sPreregConditions{}
conditions and zero in the other \sPreregEqual{}. What content hashes add
consists only of certificates that bridge a hidden change. And no
revision-ID certificate violates G1 or G2, as Theorem~3 predicts.

\subsection{First addendum: closed bridges and the schedule}
\code{protocol-addendum.json} was written after the preregistered results were
known and is labelled post hoc. \code{addendum.py} evaluates every method on
exactly the preregistered draws (bundles, read orders, start times and read
instants), so all comparisons are paired, and it first reproduces the
preregistered methods count for count. The added methods are closed-bridge
double collects with hashes and with revision IDs; the double-collect schedule
with the center page's second read removed (same instants, certified by the
exact test); the exact revisit order with content hashes; and the revisit
order in which the page read first and last offers only a content hash,
certified either by trusting the hash or by Theorem~3(ii). (The frozen addendum
protocols were written when the paper numbered this result Theorem~2 and call
it Theorem~2(ii).)

\begin{table*}[t]
\caption{Certified runs per stratum and request gap (of 12{,}000), first addendum. Revision-ID double collect $\to$ center read once $\to$ revisit order is the schedule ablation behind the $+29$ and $+122$ of Figure~2b.}
\label{tab:addendum}
\small
\begin{tabular}{@{}llrrrrrrr@{}}
\toprule
Stratum & Gap & \multicolumn{2}{c}{Double collect} & \multicolumn{3}{c}{Revisit order} & \multicolumn{2}{c}{Hash page at the ends} \\ \cmidrule(lr){3-4}\cmidrule(lr){5-7}\cmidrule(l){8-9}
 & & hash & ID & ID, center once & ID & hash & trusted & Theorem~3 \\
\midrule
vital & 1\,s & 11,997 & 11,997 & 11,997 & 11,997 & 11,997 & 11,997 & 0 \\
vital & 1\,min & 11,929 & 11,924 & 11,927 & 11,941 & 11,945 & 11,942 & 0 \\
vital & 1\,h & 10,195 & 9,954 & 10,048 & 10,279 & 10,470 & 10,341 & 0 \\
random & 1\,s & 12,000 & 12,000 & 12,000 & 12,000 & 12,000 & 12,000 & 0 \\
random & 1\,min & 11,997 & 11,997 & 11,997 & 11,997 & 11,997 & 11,997 & 0 \\
random & 1\,h & 11,852 & 11,848 & 11,851 & 11,869 & 11,873 & 11,870 & 0 \\
\bottomrule
\end{tabular}

\end{table*}

Closed bridges add exactly the runs with touching windows: at 1\,s the closed
double collect certifies the same runs as the exact test, and the added
certificates bring no violation. Over all conditions content hashes then
issue \closedHashCertificates{} certificates and revision IDs
\closedIdCertificates{}, and under both schedules the hash surplus equals the
hash certificates that break continuity: \closedHashFalseGTwo{} with the double
collect and \sHashExactContinuity{} with the revisit order. The exact
revision-ID rule never certifies fewer runs than the closed double collect,
more in \exactMoreConditions{} of \nConditions{} conditions, mostly where pages
change at the gap. For eight vital pages an hour apart the gain from
\vitalEightIdBridge{} to \vitalEightExact{} splits into
\vitalEightCenterGain{} runs from no longer requiring the center page to stay
unchanged between two reads (same instants) and \vitalEightShorterGain{} from
the shorter schedule, whose second half starts one slot earlier. One hash page
at the ends, trusted as if sound, breaks continuity \sTrustedContinuity{}
times; Theorem~3(ii) certifies none of these runs, because neither of its
windows lies inside the other pages' common interval.

\subsection{Second addendum: what a certificate asserts}
\code{protocol-addendum-2.json} was written after the first addendum and
before \code{addendum2.py} ran; the script first reproduces all 162 rows of the
first addendum. Because the preregistered G1 is existential over the claimed
interval for double collects but evaluates the single read for the exact test,
it does not measure one assertion across methods. For every certified run the
second addendum checks three nested assertions:
\begin{itemize}
\item \emph{continuity}: every twice-read page keeps its observed content
between its first and last true reads (G2);
\item \emph{claim}: what the certificate's proof asserts on the true history,
namely that every page has its observed content throughout the closed claimed
interval $[\max e^-,\min s^+]$ of a double collect, or at the true read
instant of an exact test's single read;
\item \emph{existence}: the combination held at some instant between the first
request's start and the last request's end; a certificate that fails it is
\emph{fictitious}.
\end{itemize}
The script asserts on every run that a fictitious certificate breaks its claim
and that a broken claim implies a continuity violation. It also adds the
field-aware schedule of Corollary~4: the hash-only page of the off-center
method moves to the center, the order becomes
$\mathit{order}[1{:}]+[\mathit{order}[0]]$, and the revisit order is read at the
preregistered slots. Pages then occupy different slots than in the
all-revision-ID revisit order, so this method is paired with the others by
bundle, start time and slot instants, not by page-to-slot assignment.

\begin{table*}[t]
\caption{The three nested assertions per method: certified runs and runs that break continuity, their claim, the preregistered G1, and existence (fictitious). Left: totals over both strata, $n=2,4,8$ and the three gaps. Right: eight vital pages with request starts 1\,h apart, the condition of Figure~2b. $^{\ast}$Preregistered.}
\label{tab:assertions}
\small
\setlength{\tabcolsep}{4.5pt}
\begin{tabular}{@{}lrrrrrrrrrr@{}}
\toprule
Method & \multicolumn{5}{c}{All conditions, 72{,}000 runs} & \multicolumn{5}{c}{Eight vital pages, 1\,h, 4{,}000 runs} \\ \cmidrule(lr){2-6}\cmidrule(l){7-11}
 & Certified & Continuity & Claim & G1 & Fictitious & Certified & Continuity & Claim & G1 & Fictitious \\
\midrule
content hash, $2n$, half-open$^{\ast}$ & 45,973 & 250 & 56 & 8 & 0 & 2,707 & 168 & 33 & 7 & 0 \\
revision ID, $2n$, half-open$^{\ast}$ & 45,723 & 0 & 0 & 0 & 0 & 2,539 & 0 & 0 & 0 & 0 \\
revision ID, exact, $2n{-}1^{\ast}$ & 70,083 & 0 & 0 & 0 & 0 & 2,690 & 0 & 0 & 0 & 0 \\
content hash, $2n$, closed & 69,970 & 250 & 56 & 8 & 0 & 2,707 & 168 & 33 & 7 & 0 \\
revision ID, $2n$, closed & 69,720 & 0 & 0 & 0 & 0 & 2,539 & 0 & 0 & 0 & 0 \\
revision ID, center read once & 69,820 & 0 & 0 & 0 & 0 & 2,568 & 0 & 0 & 0 & 0 \\
content hash, exact, $2n{-}1$ & 70,282 & 199 & 12 & 12 & 0 & 2,839 & 149 & 10 & 10 & 0 \\
hash page at the ends, trusted & 70,147 & 64 & 1 & 1 & 0 & 2,726 & 36 & 1 & 1 & 0 \\
hash page at the ends, Theorem~3 & 0 & 0 & 0 & 0 & 0 & 0 & 0 & 0 & 0 & 0 \\
hash page at the center, Theorem~3 & 70,032 & 0 & 0 & 0 & 0 & 2,653 & 0 & 0 & 0 & 0 \\
\bottomrule
\end{tabular}

\end{table*}

\begin{table}[t]
\caption{The hash-only page read once at the center, against the all-revision-ID revisit order on the same bundles and start times (certified of 4{,}000; the first method breaks no assertion anywhere).}
\label{tab:center}
\small
\begin{tabular}{@{}lrrrr@{}}
\toprule
Stratum & $n$ & 1\,s & 1\,min & 1\,h \\
\midrule
vital & 2 & 4,000 / 4,000 & 3,998 / 3,999 & 3,936 / 3,930 \\
vital & 4 & 4,000 / 4,000 & 3,990 / 3,989 & 3,639 / 3,659 \\
vital & 8 & 3,999 / 3,997 & 3,948 / 3,953 & 2,653 / 2,690 \\
random & 2 & 4,000 / 4,000 & 4,000 / 4,000 & 3,998 / 3,996 \\
random & 4 & 4,000 / 4,000 & 4,000 / 4,000 & 3,980 / 3,980 \\
random & 8 & 4,000 / 4,000 & 3,997 / 3,997 & 3,894 / 3,893 \\
\bottomrule
\end{tabular}

\end{table}

No certificate of any method is fictitious within the sampled runs and the
reconstructed histories (Table~\ref{tab:assertions}); the captured combination
occurs somewhere in the capture period. The median departures quoted above do
not by themselves explain this zero, nor do they imply that fictitious
certificates cannot occur: the MinIO execution of Figure~1a is one, since its
combination held at no instant. Of the \closedHashFalseGTwo{} hash
certificates that broke continuity in the closed double collect,
\closedHashClaim{} asserted an interval in which some page had other content
and \closedHashFalseGOne{} an interval in which the combination never held,
although it held elsewhere in the capture. Revision IDs break no assertion
under any schedule, and moving the hash-only page to the center restores
\mixedCenterCertificates{} of \nBundleRuns{} runs to certification, against 0
when the same page is read first and last (Table~\ref{tab:center}).

\subsection{Raw-response binding}
\code{derive.py} checks the SHA-256 of every raw API response, rebuilds every
history from the raw history responses and compares page identities and
revisions with the manifest, and checks the predecessor revisions. The request
and continuation chains are checked at collection time instead:
\code{collect.py --resume} replays the saved responses through the same request
code and stops at any divergence. For a page without a revision inside
the window, the initial state comes from a separate raw response (38 such
responses). An earlier loader checked that response's hash but took the
initial state from the manifest's copy without comparing the two, so editing
the manifest alone could change a page's initial state while every raw file
stayed intact; with the old loader such an edit to the record of
\emph{Frederick Goldie} (\code{raw/00324-before-150.json}) was accepted. The
loader now binds that state to its raw response, checking page identity,
response index, cardinality and coverage, and rejects any difference;
\code{tests/test_revision_bundle_before_window.py} keeps a valid fixture and
rejects a manifest-only edit. After the fix all eight revision-analysis outputs
are byte-identical to the archived ones, and the analysis copy in
\code{figure-artifacts/06-revision-histories/} is the same file. No reported
count changed.

\subsection{Scope}
The frozen responses give visible revision metadata and SHA-1 values, not
rendered HTML, templates, caches, serving times or administrative logs, and
timestamps have one-second resolution. The vital frame is curated and the
random frame follows \code{page_random}; neither samples the Web, and the shares
describe the retained histories of these frames. Revision IDs advance along
the reconstructed timeline; administrative deletion and undeletion, which can
re-expose a preserved revision ID after absence or after another revision
(Section~\ref{sec:platforms}), are outside the replay. Their absence from the
reconstruction is a limit of scope, not evidence that they did not occur.
